% Mục 3.3. Khối C. Người viết: Khang. Chi tiết huấn luyện nằm ở Mục 4.3.
Mỗi vùng chữ được cắt theo khung bao, đệm thêm một chút viền, rồi đưa vào VietOCR~\cite{vietocr} kiến trúc \texttt{vgg\_transformer}.
Bộ mã hóa VGG19-BN~\cite{simonyan2015vgg} biến ảnh dòng chữ thành chuỗi đặc trưng, và bộ giải mã Transformer~\cite{vaswani2017attention} sinh từng ký tự trong bảng chữ tiếng Việt có dấu.
Nhóm giải mã tham lam, không dùng beam search, để giữ tốc độ.

Phương án chính C3 là tinh chỉnh VietOCR từ trọng số đã công bố trên chính các vùng chữ đơn thuốc, gồm cả lớp \texttt{other} để mô hình gặp đủ các kiểu chữ.
Nhóm cũng nâng chiều rộng tối đa của ảnh vào từ 512 lên 1024~px, vì với 512 những vùng chữ dài như dòng chẩn đoán bị thu nhỏ theo chiều ngang và mất chữ \chuong{2}.

Có ba phương án đối chứng.
C2 là VietOCR với trọng số đã công bố, không tinh chỉnh.
C0 và C1 dùng Tesseract 5~\cite{smith2007tesseract} với gói ngôn ngữ tiếng Việt, đọc từng vùng như một dòng chữ.
C0 nhận ảnh xám phóng lớn hai lần bằng nội suy bậc ba, tức thay đổi độ phân giải không gian của vùng cắt \chuong{2}, còn C1 giống C0 nhưng đầu vào là ảnh nhị phân Sauvola \chuong{5}, nhằm kiểm tra xem nhị phân hóa có giúp nhận dạng hay không.
